多模态情感预测面临三重困难：文本、语音与视觉在采样频率与时间尺度上差异显著，实际采集又常出现局部时段模态不可用，而多数模型的决策依据难以复核。本文以赛题提供的CMU-MOSEI系列数据为唯一数据来源，围绕时序特征构建、缺失鲁棒预测与可解释分析，依次建立不确定边界锚点对齐模型、教师--学生鲁棒预测模型与模态子集训练及输入干预解释模型。

针对问题一，以原词为时间锚点、以子词为输出位置，结合BERT语义表示、25维声学描述符与16维人脸几何及运动描述符，采用CTC强制对齐与时间交叠池化建立三模态对应关系，并由观测覆盖、路径置信度与边界扰动合成锚点质量分数。完成附件1全部100条样本的处理：三模态特征维度分别为$50\times768$、$50\times25$与$50\times16$；1932个原词中1920个生成时间区间，区间生成率为99.38\%。全量清单、词级明细与帧级记录支持从特征、词区间到原始帧段的逐级追溯，典型样本给出文字、语音与视频的对应关系。

针对问题二，建立三态感知可靠补全与融合模型TSRF-D。以内容掩码与观测掩码区分填充、观测与局部缺失三种状态，通过观测节点加权补全、可靠度门控融合与双向时序编码联合预测情感极性与连续强度，训练阶段由真实标签监督、教师任务蒸馏与重建辅助目标共同约束。验证集Macro-F1为0.6012、强度MAE为0.5915，附件2测试集分别为0.6138与0.6275；8锚点局部文本缺失时F1降至0.5198，而局部语音、视觉缺失时分别为0.6157与0.6100，表明文本证据对当前模型的影响最大。补全相对回退使特征MAE降低0.0260，改善主要来自视觉模态，蒸馏与门控在多种缺失条件下给出可核验的对照结果。模型完成附件3全部30条样本的极性与强度预测。

针对问题三，构建ICF-Shapley解释模型，以模态子集训练覆盖解释输入状态，在固定参考类别下精确枚举模态联盟，并通过局部输入删除定位证据片段。验证集Macro-F1为0.6170、MAE为0.5954，附件2测试集分别为0.6054与0.6303；96条验证样本中，高分片段删除相对随机删除的AOPC差为0.1068，95\%区间为[0.0887,0.1259]。完成附件4全部20条样本的预测与解释，结合词级CTC区间与真实视频PTS输出文本、语音与关键帧证据；不同随机种子与不同解释基线下的主要参考模态一致率分别达到82\%以上与87\%以上，模态贡献具有较高一致性。

\keywords{多模态情感\quad 时序对齐\quad 局部模态缺失\quad 教师--学生蒸馏\quad Shapley解释}
